\documentclass[11pt]{article}

\usepackage{amsmath,amssymb,amstext,amsthm}
\usepackage{geometry}
\usepackage{fancyhdr}
\usepackage[pdftex]{graphicx}
\usepackage{xcolor}

\geometry{
  verbose,
  dvips,
  width=430pt, marginparsep=0pt, marginparwidth=0pt,
  top=90pt, headheight=12pt, headsep=20pt, footskip=30pt, bottom=80pt}

\setlength{\parskip}{\medskipamount}
\setlength{\parindent}{0pt}

\linespread{1.05}

\newtheorem{theorem}{Theorem}
\newtheorem{lemma}[theorem]{Lemma}
\newtheorem{proposition}[theorem]{Proposition}
\newtheorem{corollary}[theorem]{Corollary}
\newtheorem{conjecture}[theorem]{Conjecture}
\newtheorem{obs}{Observation}
\newtheorem{clm}{Claim}
\newtheorem{ques}{Question}
\newtheorem{problem}{Problem}

\newtheorem{ex}{Example}


\newcommand{\spc}{\hspace{0.08in}}
\newcommand{\vs}{\vspace*{.1in}}
\newcommand{\E}{\mathbb{E}}
\newcommand{\Prob}{\mathbb{P}}
\newcommand{\edit}[1]{\emph{\textcolor{blue}{#1}}}
\newcommand*\dif{\mathop{}\!\mathrm{d}}


\renewenvironment{proof}{{\noindent \textbf{\textit{Proof.}}}}
{\hfill $\Box$ \vspace*{0.1in}}
\newenvironment{proofc}{{\noindent \textit{Proof of Claim.}}}
{\hfill $\Box$}



\begin{document}

\begin{LARGE}\begin{center}\bf 16.1 Vector Fields\end{center}
\end{LARGE}

\vs\vs



\section{Warm Up}
\textbf{Problem 1:} Find the gradient of $f(x,y) = x^2 + e^{xy}$.

\section{Motivation}
We switch our attention to a new definition called vector fields. As of right now we take the definition as is, as we develop some intuition for what these are and how they could be useful!



\section{Definitions \& Theorems}

\begin{enumerate}
\item Our story of vector fields, of course begins with the definition of one. A vector field on some set  $D\subseteq \mathbb{R}^2$ is a function $F$ which assigns to each point $(x,y)$ a 2-dimensional vector $F(x,y)$.

\item A vector field on some set $E\subseteq \mathbb{R}^3$ is a function $F$ which assigns each point $(x,y,z)$ a 3-dimensional vector $F(x,y,z)$.


\item If we remember back from Chapter 14, the gradient $\nabla f(x,y) = \langle f_x(x,y), f_y(x,y)\rangle$. Hence the gradient is really a vector field over $\mathbb{R}^2$ and hence we call it a gradient vector field. If $F(x,y)$ is the gradient of some function $f(x,y)$, we call $F$ a conservative vector field.

\item For right now our sole objective is to get some intuition towards what vector fields look like and how we can visualize them. 
\end{enumerate}




\section{Examples}

\begin{problem}
    Draw the following vector fields

    \begin{enumerate}
        \item $F(x,y) = 2\hat{i}$
        \item $F(x,y) = -x\hat{i} - y\hat{j}$
        \item $F(x,y) = y\hat{i} - x\hat{j}$
        \item $F(x,y) = x^2 \hat{i} + y^2\hat{j}$
        \item $F(x,y) = y^2 \hat{i} + x^2 \hat{j}$
        \item $F(x,y) = \frac{x}{\sqrt{x^2+y^2}}\hat{i} + \frac{y}{\sqrt{x^2+y^2}}\hat{j}$
        \item $F(x,y) = \hat{i} + x\hat{j}$
        \item $F(x,y) = x \hat{i} + \hat{j}$
    \end{enumerate}
\end{problem}



\section{Scratch Work}

\end{document} 